\section{Introduction}

Many of society’s most time-critical predictions—from tsunami early warnings to flood inundation, structural failure, and subsurface pollutant transport alerts—are underpinned by high-fidelity partial differential equation (PDE) solvers that capture fine-grained physics \citep{leveque2002finite,palais2009differential, quarteroni2010numerical}. However, their computational cost poses a persistent challenge in time-critical or real-time settings. For example, warning systems for tsunamis, floods, volcanic plumes, or landslides have only minutes to digest real-time sensor data and produce actionable forecasts, whereas a single high-resolution simulation may take hours. Inversions or learning tasks require thousands of runs \citep{li2021fourier} with high-dimensional parametric inputs such as geometries, bed topographies, and roughnesses \citep{grady2023model, xiao2024fourier}. Since real-time tsunami modeling is unfeasible \citep{Reymond2012}, current warning systems precompute and store many scenarios and, during an event, superimpose the closest ones to generate predictions \citep{gica2008development, fujita2024scenario}.

Dynamic source modeling approaches that embed fault slip directly within high-resolution elasticity solvers have also been proposed to capture transient rupture effects prior to steady-state deformation \citep{vogl2017high}. However, such approaches remain computationally intensive and do not eliminate the fundamental cost barrier.

Consequently, reduced-fidelity approximations can introduce substantial errors and fail to capture critical system dynamics. For example, early tsunami forecasts during the 2011 Tōhoku event substantially underestimated wave heights, limiting the effectiveness of initial warnings \citep{hoshiba2014earthquake, ozaki2012jma}. Similar challenges arise in weather forecasting and system controls. This tension motivates a central question: can high-fidelity physical modeling be made practical for real-time forward and inverse applications?

Strong hopes were ignited recently by artificial intelligence for scientific computing and neural network surrogates \citep{zong2026mathematics}, such as neural operators (NOs), which promise many orders of magnitude in acceleration for PDE surrogates and real-time, super-resolution inference \citep{li2021fourier, lu2021learning, kovachki2023neural, wang2024deep, qin2024toward}. NOs, e.g., Fourier Neural Operators (FNO) \citep{li2021fourier}, Wavelet Neural Operators (WNO) \citep{tripura2023wavelet}, DeepONets \citep{lu2021learning}, and Convolutional Neural Operators (CNO) \citep{raonic2023convolutional}, learn efficient mappings between function spaces from numerical simulations and have seen applications in lithography (chip design) \citep{yang2022large}, weather forecasting \citep{kurth2023fourcastnet}, fluid mechanics \citep{han2022equivariant}, and subsurface CO$_2$ sequestration \citep{wen2023real}. However, NOs still face an existential limitation (``curse of dimensionality") \citep{lanthaler2025parametric, kovachki2024data}: purely empirical training struggles to represent high-dimensional or spatially distributed inputs such as heterogeneous parameters, geometries, and forcing fields, leading to large training data demands and challenges with inversion for nonlinear equations \citep{kovachki2023neural}. This unfavorable scaling means NOs degrade sharply and need to be retrained when gridded inputs change noticeably, posing a central barrier to their deployment. We therefore ask whether it is possible to mitigate the data burden associated with large input dimensionality (``the curse of dimensionality") while retaining high-fidelity surrogate accuracy.

A key but underappreciated failure mode of neural operators is not inaccurate predictions, but learning inaccurate sensitivities—how solutions change with respect to inputs—leading to poor performance under input perturbations, out-of-distribution shifts, and inverse settings \citep{choi2024applications, behroozi2025sensitivity}. Several approaches have sought to incorporate physical constraints into surrogate training. Enforcing local PDE loss \citep{li2024physics} or equivariance \citep{han2022equivariant} covers some possible directions. Another group of methods enforces consistency between integrated sensitivities through the NOs and those from numerical solvers, including Sensitivity-Constrained Fourier Neural Operators (SC-FNO) \citep{behroozi2025sensitivity}, derivative-enhanced DeepONet (DE-DeepONet) \citep{qiu2024derivative}, derivative-informed neural operators (DINO) \citep{cao2025derivative}, and Sobolev training \citep{czarnecki2017sobolev}. These methods improve generalization, but to date have largely been limited to low-dimensional parametric settings. Many rely on compressed or low-rank representations involving only tens of parameters, as full Jacobian supervision is widely regarded as computationally prohibitive. These approaches address a complementary representation bottleneck and are, in principle, compatible with sensitivity supervision proposed below. In surrogate training settings, an important insight is that the size of the sensitivity Jacobian is proportional to the degrees of freedom (DOF). We seek to leverage this insight and ask whether we can constrain high-dimensional gridded inputs with DOF exceeding $10^4$.

Here, we identify a previously unrecognized ``sensitivity-driven scaling'' principle in operator learning: physical sensitivities constitute an information channel that scales in capacity with system input size and can cost-effectively reduce uncertainty. Building on this principle, we introduce the sensitivity-constrained neural operator (SC-NO), which matches physics-derived Jacobians to those of neural operators, thereby enabling full-fidelity forward and inverse solutions even in high-dimensional, time-critical settings (theoretical justification is provided in Appendix~\ref{sec_sup_method}). The cost of sampling the exact full Jacobian is amortized within the surrogate training loop. This effect of this unique scaling behavior is demonstrated in a tsunami-warning case study involving real-time inversion of a gridded seafloor deformation field from sparse buoy observations (Figure~\ref{fig:stage3_param_inversion}), followed by accurate forecasting of coastal impacts—all completed within 4.5 minutes on a single graphical processing unit (GPU). We then employ diverse controlled PDE experiments to show how standard neural operator challenges are related to their data-scaling behavior, while sensitivity supervision fundamentally alters this behavior. Beyond accuracy, this constraint improves out-of-distribution generalization and long-horizon stability, which are critical for inversion. The central contribution is not a new neural architecture, but the discovery of how sensitivity supervision fundamentally alters surrogate scaling laws and enables previously implausible high-dimensional applications.
